\documentclass[11pt]{article}

\usepackage[T1]{fontenc}
\usepackage[utf8]{inputenc}
\usepackage{amsmath,amssymb,amsthm}
\usepackage{array,booktabs}
\usepackage[margin=1in]{geometry}
\usepackage[hidelinks]{hyperref}

\newtheorem{theorem}{Theorem}
\newtheorem{lemma}[theorem]{Lemma}
\newtheorem{proposition}[theorem]{Proposition}
\newtheorem{corollary}[theorem]{Corollary}
\theoremstyle{definition}
\newtheorem{definition}{Definition}
\newtheorem{remark}{Remark}

\newcommand{\leanus}{\textunderscore\allowbreak}
\newcommand{\lean}[1]{\ifmmode\text{\ttfamily #1}\else{\ttfamily\let\_\leanus #1}\fi}
\newcommand{\TC}{\operatorname{TC}}
\newcommand{\bTC}{\underline{\operatorname{TC}}}
\newcommand{\PT}{\operatorname{PT\text{-}rank}}
\newcommand{\SoS}{\operatorname{SoS}}
\newcommand{\Prod}{\operatorname{Prod}}
\newcommand{\rk}{\operatorname{rank}}
\newcommand{\tr}{\operatorname{tr}}
\newcommand{\supp}{\operatorname{supp}}
\newcommand{\sym}{\operatorname{sym}_4}
\newcommand{\K}{\mathbb{K}}
\newcommand{\R}{\mathbb{R}}
\newcommand{\Q}{\mathbb{Q}}
\newcommand{\Z}{\mathbb{Z}}
\newcommand{\C}{\mathbb{C}}

\title{Two-cut rank, partial-transpose rank, and sum-of-squares complexity of biquadratic forms}
\author{PastaPastaPasta}
\date{November 2026}

\emergencystretch=2em

\begin{document}
\maketitle

\begin{abstract}
We study the \emph{two-cut rank} $\TC_\K(M) := \min\{\rk A + \rk W : M = A + W^\Gamma\}$ of an $rs \times rs$ matrix $M$ under the partial transpose $(\cdot)^\Gamma$, and its relation to the \emph{partial-transpose rank} $\PT_\K(M)$ of Rossman and Zhu~\cite{RZ25}, the bilinear product complexity $\Prod_\K(Q_M)$, and the sum-of-squares complexity $\SoS_\K(Q_M)$ of the associated biquadratic form $Q_M(x, y) = \sum M_{(a,b),(c,d)} x_a x_c y_b y_d$. Every formalised theorem and counterexample in this paper is unconditional and kernel-checked in Lean~4 with Mathlib (library \lean{PastaMath.TwoCutSoS}); the single paper-only bridge is noted explicitly below and in Section~\ref{sec:verified}.

\begin{enumerate}
\item \textbf{Exact dictionary and refutation of Rossman--Zhu \S3.2.} Over every field $\K$, $\PT_\K(M) = \TC_\K(M)$ and $\Prod_\K(Q_M) \le \TC_\K(M)$. When $\operatorname{char}\K \neq 2$ and $M$ is fully symmetric, $\TC_\K(M) \le 2\SoS_\K(Q_M)$ and $\TC_\K(M) \le 4\Prod_\K(Q_M)$. When $\operatorname{char}\K \neq 2$ and $\sqrt{-1} \in \K$ (for arbitrary $M$), $\SoS_\K(Q_M) \le 2\Prod_\K(Q_M) \le 2\TC_\K(M) = 2\PT_\K(M)$. We exhibit explicit matrices ($M_1, M_3 \in M_4(\K)$ and $M_2 \in M_{16}(\K)$) refuting the printed statements of \cite[Lemma~3.13]{RZ25} ($\SoS(Q_M) \le \PT(M)$) and both inequalities of \cite[Proposition~3.16]{RZ25} ($\SoS(Q_M)-1 \le \PT(M) \le 2^{d-1}\SoS(Q_M)$) at $d=2$ over every field of characteristic $\neq 2$, and show that the repaired factor $2$ in $\SoS_\K(Q_M) \le 2\PT_\K(M)$ is tight for general $M$. Whether $\SoS_\K(Q_M) \le \PT_\K(M)$ holds for \emph{fully symmetric} $M$ over algebraically closed fields remains open.
\item \textbf{Characteristic-free crossing bounds.} For the identity matrix $I_{rs}$ (corresponding to the Hurwitz biquadratic form $X_{r,s} = (\sum_{i=1}^r x_i^2)(\sum_{j=1}^s y_j^2)$), over \emph{every} field $\K$ including $\operatorname{char}\K = 2$: $\TC_\K(r,s) \ge \max(r,s)$ for all $r,s \ge 1$; $\TC_\K(r,s) \ge \max(r,s)+1$ whenever $r,s \ge 3$; $\TC_\K(2,s) = 2\lceil s/2\rceil$; $\TC_\K(r,s) \le 1 + rs - \binom{\min(r,s)+1}{2}$; and hence $\TC_\K(3,3) = 4$.
\item \textbf{Explicit exact and border witnesses.} Over every field $\K$, explicit $\pm 1$ decompositions give $\TC_\K(4,5) \le 8$, $\TC_\K(4,6) \le 8$, and $\TC_\K(8,8) \le 16$, and a pole-order-$2$ border decomposition gives $\bTC_{\K,2}(4,4) \le 6$. Over $\operatorname{char}\K \neq 2$, $n$ anticommuting involutions in $M_d(\K)$ yield $\TC_\K(d,n) \le 2d$; over every field $\K$, a pole-order-$k$ border two-cut decomposition of $I_{rs}$ of ranks $a + b$ yields $\TC_\K(r,s) \le (k+1)(a+b)$.
\item \textbf{Word-span tensor powers of reflection-type decompositions.} For any reflection-type decomposition ($\operatorname{spec} M = \{0^{D-n},\lambda^n\}$, $\operatorname{spec} M^\Gamma = \{1^{D-n},(1-\lambda)^n\}$, $D = rs$, $p = n/D$, $\lambda = 1/(2p)$), word-span decompositions of $K$-th tensor powers reduce to a weighted uncertainty problem for the $p$-biased reflection operator $T_{-1} = E_1(p)^{\otimes K}$ in $L^2(\mu_p)$. When $\lambda = a/b$ with a prime $\ell \mid (a-b)$ and $\ell \nmid a$, we prove a \emph{Valuation Theorem}: for every $K$, no $\mathbb{Z}_{(\ell)}$-linear word combination beats the Bayes error $\sum_T \min(\mu_p(T), \mu_{1-p}(T))$, attained exactly by the layer grouping. Finally, via tensorising Farkas-certificate packings---parametric at $K = 1, 2$ together with a kernel-checked $22$-certificate rational packing of exact value $57/256$ at $K = 4$, $p = 1/4$---we prove that for every shape $r,s \ge 2$ and every reflection-type decomposition over $\R$, no word-span tensor power achieves asymptotic rate $\le \sqrt{1+q^2}$ ($q = \log\min(r,s)/\log\max(r,s)$).
\end{enumerate}

\emph{What is kernel-checked:} every formalised theorem and counterexample listed in Table~\ref{tab:lean}, with zero \lean{sorry}, zero \lean{native\_decide}, and only the standard axioms \lean{[propext, Classical.choice, Quot.sound]}.
\emph{What is not machine-checked:} the Kronecker-product simultaneous-diagonalisation bridge identifying the matrix rank of a $K$-th word-span tensor power with $D^K \varphi_K(p; h)$ is proved on paper; furthermore, the word-span no-go theorem applies to word-span decompositions of reflection-type diagonalisable bases over $\R$ and does not rule out general (non-word-span) decompositions, non-reflection spectra, non-diagonalisable bases, or border families.
\end{abstract}

\section{Introduction}

\paragraph{Background.}
For a biquadratic form $Q(x,y)$ in variables $x = (x_1,\dots,x_r)$ and $y = (y_1,\dots,y_s)$ over a field $\K$, its \emph{sum-of-squares (SoS) complexity} $\SoS_\K(Q)$ is the minimum $c$ such that $Q(x,y) = \sum_{k=1}^c L_k(x,y)^2$ with each $L_k$ bilinear, and its \emph{bilinear product complexity} $\Prod_\K(Q)$ is the minimum $c$ such that $Q(x,y) = \sum_{k=1}^c U_k(x,y) V_k(x,y)$ with $U_k, V_k$ bilinear. Hrube\v{s}, Wigderson and Yehudayoff~\cite{HWY11} proved that a superlinear lower bound $\SoS_\C(Q) = n^{1+\Omega(1)}$ for an explicit biquadratic form implies exponential lower bounds for non-commutative arithmetic circuits. Rossman and Zhu~\cite{RZ25} extended this connection to $d$-multiquadratic forms and $\mathrm{VNC}^1 \neq \mathrm{VNP}$, introducing the \emph{partial-transpose rank} $\PT_\K(M)$ of the coefficient matrix $M$ representing $Q$.

For the canonical Hurwitz form $X_{r,s}(x,y) = (\sum_{i=1}^r x_i^2)(\sum_{j=1}^s y_j^2)$ (represented by $M = I_{rs}$), classical Hurwitz--Radon composition gives $\SoS_\R(X_{r,s}) = r$ only for Hurwitz--Radon pairs, while over $\C$ the asymptotic exponent $\lim_{N\to\infty} \log_N \SoS_\C(X_{N,N}) \le \sqrt{2}$ was recently established by Siemaszko~\cite{Sie26} (we make no priority claim on that square exponent). A natural programme for understanding both lower bounds and rectangular exponents $\tau(q)$ (where $r = N$, $s = N^q$, $0 < q \le 1$, and the square-derived benchmark is $\sqrt{1+q^2}$) is to work directly with the matrix \emph{two-cut rank}
\[
  \TC_\K(M) := \min\bigl\{\rk A + \rk W : A, W \in \K^{(r\times s)\times(r\times s)},\; M = A + W^\Gamma\bigr\},
\]
where $(W^\Gamma)_{(a,b),(c,d)} := W_{(a,d),(c,b)}$ is the partial transpose on the second index.

\paragraph{Overview of results.}
All theorem names in typewriter font refer to the Lean~4 library \lean{PastaMath.TwoCutSoS}; Section~\ref{sec:verified} gives the complete theorem-to-Lean index.

\begin{itemize}
\item \textbf{Section~\ref{sec:dict} (Dictionary and Rossman--Zhu \S3.2).} We prove $\PT_\K(M) = \TC_\K(M)$ (\lean{hasPTRank\_iff\_hasTwoCut}) and $\Prod_\K(Q_M) \le \TC_\K(M)$ (\lean{prod\_of\_twoCut}) over every field $\K$; $\TC_\K(M) \le 2\SoS_\K(Q_M)$ and $\TC_\K(M) \le 4\Prod_\K(Q_M)$ for fully symmetric $M$ when $\operatorname{char}\K \neq 2$ (\lean{twoCut\_of\_sos}, \lean{twoCut\_of\_prod}); and $\SoS_\K(Q_M) \le 2\Prod_\K(Q_M) \le 2\TC_\K(M) = 2\PT_\K(M)$ when $\operatorname{char}\K \neq 2$ and $\sqrt{-1} \in \K$ (\lean{sos\_of\_prod}, \lean{sos\_of\_twoCut}, \lean{sos\_of\_ptRank}). We then construct three explicit matrices ($M_1, M_3 \in M_4(\K)$ and $M_2 \in M_{16}(\K)$) refuting the printed statements of \cite[Lemma~3.13]{RZ25} and both sides of \cite[Proposition~3.16]{RZ25} at $d = 2$ over every field of characteristic $\neq 2$ (\lean{rz\_lemma313\_counterexample}, \lean{rz\_prop316\_left\_counterexample}, \lean{rz\_prop316\_right\_counterexample}), showing that the factor $2$ in $\SoS_\K(Q_M) \le 2\PT_\K(M)$ is sharp in general.
\item \textbf{Section~\ref{sec:lb} (Characteristic-free crossing bounds).} Over \emph{every} field $\K$ (including $\operatorname{char}\K = 2$), we prove $\TC_\K(r,s) \ge \max(r,s)$ for all $r,s \ge 1$ (\lean{twoCut\_ge\_max}), $\TC_\K(r,s) \ge \max(r,s)+1$ for all $r,s \ge 3$ (\lean{twoCut\_gt\_max\_of\_ge\_three}), the exact formula $\TC_\K(2,s) = 2\lceil s/2\rceil$ (\lean{twoCut\_two\_iff}), the transpose upper bound $\TC_\K(r,s) \le 1 + rs - \binom{\min(r,s)+1}{2}$ (\lean{transpose\_twoCut}), and $\TC_\K(3,3) = 4$ (\lean{twoCut\_three\_three\_iff}).
\item \textbf{Section~\ref{sec:witnesses} (Explicit exact and border witnesses).} We formalise the self-dual Clifford bound $\TC_\K(d,n) \le 2d$ (\lean{clifford\_twoCut}), border-to-exact conversion $\TC_\K(r,s) \le (k+1)(a+b)$ (\lean{exact\_of\_border\_coeffs}), explicit $\pm 1$ integer schemes proving $\TC_\K(4,5) \le 8$, $\TC_\K(4,6) \le 8$, and $\TC_\K(8,8) \le 16$ over every field (\lean{tc45\_twoCut}, \lean{tc46\_twoCut}, \lean{tc88\_twoCut}), and an explicit pole-order-$2$ border decomposition proving $\bTC_{\K,2}(4,4) \le 6$ over every field (\lean{border44\_rank\_le}, \lean{border44\_laurent\_identity}, \lean{border44\_identity\_field}).
\item \textbf{Section~\ref{sec:refl} (Word-span tensor powers of reflection-type decompositions).} A major class of exact two-cut decompositions---including Hurwitz $[2,2,2]$, $[4,4,4]$, $[8,8,8]$, all Clifford self-dual decompositions, the $(2,2)$ transpose point, and the $(4,6)$ scheme---have \emph{reflection-type} spectra $\operatorname{spec} M = \{0^{D-n}, \lambda^n\}$ and $\operatorname{spec} M^\Gamma = \{1^{D-n}, (1-\lambda)^n\}$. We prove that word-span decompositions of their $K$-th tensor powers are governed by the $p$-biased reflection operator $T_{-1}$ in $L^2(\mu_p)$ (\lean{reflOp\_involutive}, \lean{reflOp\_selfadjoint}, \lean{wordPoly\_eval\_B}), prove the Valuation Theorem (\lean{valuation\_theorem}, \lean{layer\_grouping\_le}) showing that layer groupings are exactly optimal at every $K$ whenever $\lambda = a/b$ has a prime $\ell \mid (a-b)$, $\ell \nmid a$, develop tensorising Farkas-certificate packings (\lean{CertPacking.lower\_bound}, \lean{CertPacking.asymptotic\_bound}), verify a $22$-certificate rational packing of value $57/256$ at $K = 4, p = 1/4$ (\lean{packing\_K4\_quarter\_value}), and prove that no reflection-type word-span tensor power beats the rate $\sqrt{1+q^2}$ for any shape $r,s \ge 2$ (\lean{reflection\_wordSpan\_nogo}).
\end{itemize}

\paragraph{AI assistance.}
The proofs and the Lean formalisation were produced with AI coding agents under the author's direction. The kernel-checked Lean proofs are the primary evidence for every theorem stated below; this paper presents the mathematical proofs and documents their exact scope.

\section{Definitions, dictionary, and Rossman--Zhu \S3.2}\label{sec:dict}

\subsection{Matrix partial transpose and biquadratic forms}

Let $\K$ be a field and $r, s \ge 1$. Matrices in $M_{rs}(\K)$ are indexed by pairs $(a,b) \in [r]\times[s]$. The \emph{partial transpose} $M^\Gamma$ of $M \in M_{rs}(\K)$ is defined by
\[
  (M^\Gamma)_{(a,b),(c,d)} := M_{(a,d),(c,b)}.
\]
It is a $\K$-linear involution: $(M^\Gamma)^\Gamma = M$, $I^\Gamma = I$, and $\tr(M^\Gamma) = \tr(M)$ (\lean{ptr\_ptr}, \lean{ptr\_one}, \lean{trace\_ptr}). A matrix $M$ is \emph{fully symmetric} (\lean{IsFullySym}) if $M^\top = M$ and $M^\Gamma = M$.

The biquadratic form associated with $M$ is
\[
  Q_M(x, y) := \sum_{a,c \in [r]} \sum_{b,d \in [s]} M_{(a,b),(c,d)}\, x_a x_c y_b y_d.
\]
The coefficient of the monomial $x_a x_c y_b y_d$ in $Q_M$ is the sum of the distinct matrix entries contributing to that monomial (\lean{biquadCoeff}): $M_{(a,b),(a,b)}$ if $a=c$ and $b=d$; $M_{(a,b),(a,d)} + M_{(a,d),(a,b)}$ if $a=c$, $b\neq d$; $M_{(a,b),(c,b)} + M_{(c,b),(a,b)}$ if $a\neq c$, $b=d$; and $M_{(a,b),(c,d)} + M_{(c,b),(a,d)} + M_{(a,d),(c,b)} + M_{(c,d),(a,b)}$ if $a\neq c$ and $b\neq d$. Two matrices $M, N$ have the same formal biquadratic coefficients (\lean{SameBiquad}) if and only if $Q_M(x,y) = Q_N(x,y)$ for all $x, y$---this holds over \emph{every} field $\K$, including $\operatorname{char}\K = 2$ (\lean{sameBiquad\_iff\_qf}). When $\operatorname{char}\K \neq 2$, this is also equivalent to equality of the four-way symmetrisations
\[
  (\sym M)_{(a,b),(c,d)} := M_{(a,b),(c,d)} + M_{(c,b),(a,d)} + M_{(a,d),(c,b)} + M_{(c,d),(a,b)}
\]
(\lean{sameBiquad\_iff\_sym4}). In every characteristic, partial transpose preserves all monomial coefficients: $Q_{M^\Gamma} = Q_M$ (\lean{biquadCoeff\_ptr}).

\begin{definition}[Two-cut, PT-rank, product, and SoS predicates]\label{def:ranks}
For $M \in M_{rs}(\K)$ and $c \in \mathbb{N}$:
\begin{enumerate}
\item \emph{Two-cut rank} (\lean{HasTwoCut}): $\TC_\K(M) \le c$ if there exist $A, W \in M_{rs}(\K)$ with $M = A + W^\Gamma$ and $\rk A + \rk W \le c$. Write $\TC_\K(r,s) := \TC_\K(I_{rs})$.
\item \emph{Partial-transpose rank} (\lean{HasPTRank}, following \cite[\S3.2]{RZ25} for $d=2$): a matrix $B$ is \emph{PT-basic} (\lean{IsPTBasic}) if either $\rk B = 1$ or $\rk(B^\Gamma) = 1$; $\PT_\K(M) \le c$ if $M = \sum_{i=1}^m B_i$ with $m \le c$ and each $B_i$ PT-basic.
\item \emph{Bilinear product complexity} (\lean{HasProdFormula}): $\Prod_\K(Q_M) \le c$ if there exist $U_1, V_1, \dots, U_c, V_c \in \K^{r\times s}$ such that $\sum_{k=1}^c u_k v_k^\top$ and $M$ have the same formal biquadratic coefficients (\lean{SameBiquad}, where $u_k, v_k$ are the vectorisations of $U_k, V_k$). Equivalently, over every field $\K$, $Q_M(x,y) = \sum_{k=1}^c (x^\top U_k y)(x^\top V_k y)$ for all $x,y$ (\lean{hasProdFormula\_iff\_qf}).
\item \emph{Sum-of-squares complexity} (\lean{HasSoSFormula}): $\SoS_\K(Q_M) \le c$ if there exist $L_1, \dots, L_c \in \K^{r\times s}$ such that $\sum_{k=1}^c \ell_k \ell_k^\top$ and $M$ have the same formal biquadratic coefficients (\lean{SameBiquad}). Equivalently, over every field $\K$, $Q_M(x,y) = \sum_{k=1}^c (x^\top L_k y)^2$ for all $x,y$ (\lean{hasSoSFormula\_iff\_qf}).
\end{enumerate}
\end{definition}

\subsection{The exact dictionary}

\begin{theorem}[Dictionary; \lean{Dictionary.lean}]\label{thm:dict}
Let $\K$ be any field and $M \in M_{rs}(\K)$.
\begin{enumerate}
\item \textbf{PT-rank equals two-cut rank} (\lean{hasPTRank\_iff\_hasTwoCut}): $\PT_\K(M) \le c \iff \TC_\K(M) \le c$.
\item \textbf{Two-cut implies product} (\lean{prod\_of\_twoCut}): $\TC_\K(M) \le c \implies \Prod_\K(Q_M) \le c$.
\item \textbf{SoS implies two-cut for fully symmetric $M$} (\lean{twoCut\_of\_sos}, \cite[Lemma~3.15]{RZ25} for $d=2$): if $\operatorname{char}\K \neq 2$ and $M$ is fully symmetric, then $\SoS_\K(Q_M) \le c \implies \TC_\K(M) \le 2c$.
\item \textbf{Product implies two-cut for fully symmetric $M$} (\lean{twoCut\_of\_prod}): if $\operatorname{char}\K \neq 2$ and $M$ is fully symmetric, then $\Prod_\K(Q_M) \le c \implies \TC_\K(M) \le 4c$.
\item \textbf{Polarisation} (\lean{sos\_of\_prod}, \lean{sos\_of\_twoCut}, \lean{sos\_of\_ptRank}): if $\operatorname{char}\K \neq 2$ and $\K$ contains a square root $i$ of $-1$, then for every $M$,
\[
  \Prod_\K(Q_M) \le c \implies \SoS_\K(Q_M) \le 2c, \qquad
  \PT_\K(M) \le c \implies \SoS_\K(Q_M) \le 2c.
\]
\end{enumerate}
In particular, when $\operatorname{char}\K \neq 2$, $\sqrt{-1} \in \K$, and $M$ is fully symmetric (e.g.\ $M = I_{rs}$), $\Prod_\K(Q_M)$, $\TC_\K(M) = \PT_\K(M)$, and $\SoS_\K(Q_M)$ are equivalent up to universal constant factors: $\Prod_\K(Q_M) \le \TC_\K(M) \le 2\SoS_\K(Q_M) \le 4\Prod_\K(Q_M)$ and $\SoS_\K(Q_M) \le 2\TC_\K(M)$.
\end{theorem}

\begin{proof}[Proof sketch]
(1) Every matrix $A$ of rank $a$ is a sum of $a$ rank-$1$ matrices (\lean{exists\_sum\_rank\_one\_eq}). Grouping PT-basic summands according to whether $\rk B_i = 1$ or $\rk B_i^\Gamma = 1$ yields $M = A + W^\Gamma$ with $\rk A + \rk W \le c$, and conversely expanding $A$ and $W$ as sums of rank-$1$ matrices gives $\rk A + \rk W$ PT-basic summands.
(2) If $M = A + W^\Gamma$, write $A = \sum_{k=1}^a u_k v_k^\top$ and $W = \sum_{j=1}^w p_j q_j^\top$; since $\sym(W^\Gamma) = \sym(W)$, we have $\sym M = \sym(\sum u_k v_k^\top + \sum p_j q_j^\top)$.
(3) Given $\sym(L^\top L) = \sym M = 4M$ (since $M$ is fully symmetric) and $(L^\top L)^\top = L^\top L$, we have $2(L^\top L + (L^\top L)^\Gamma) = \sym(L^\top L) = 4M$. Setting $S := \frac{1}{2} L^\top L$ gives $S + S^\Gamma = M$ and $\rk S \le c$, so $\TC_\K(M) \le 2c$.
(4) Given $G = \sum_{k=1}^c u_k v_k^\top$ with $\sym G = 4M$, setting $S := \frac{1}{4}(G + G^\top)$ gives $S + S^\Gamma = \frac{1}{4}\sym G = M$ and $\rk S \le 2c$, so $\TC_\K(M) \le 4c$.
(5) Polarise each product via $UV = \bigl(\frac{U+V}{2}\bigr)^2 + \bigl(\frac{i(U-V)}{2}\bigr)^2$.
\end{proof}

Note that the factor $2$ in Theorem~\ref{thm:dict}(3) cannot be removed even for fully symmetric $M = I_{16}$: Euler's four-square identity gives $\SoS_\K(X_{4,4}) \le 4$, whereas Theorem~\ref{thm:lb}(2) below proves $\TC_\K(4,4) \ge 5 > 4$.

\subsection{Refutation of Rossman--Zhu Lemma~3.13 and Proposition~3.16}\label{sec:rz}

In~\cite[\S3.2]{RZ25}, Rossman and Zhu state for $d$-multiquadratic forms (specialised to $d=2$):
\begin{itemize}
\item \textbf{Lemma~3.13} (printed without a full-symmetry hypothesis on $M$): $\SoS(Q_M) \le \PT(M)$.
\item \textbf{Proposition~3.16} (printed without a full-symmetry hypothesis on $M$): $\SoS(Q_M) - 1 \le \PT(M) \le 2^{d-1}\SoS(Q_M)$.
\end{itemize}

\begin{theorem}[Counterexamples to \cite{RZ25} Lemma~3.13 and Proposition~3.16; \lean{RossmanZhu.lean}]\label{thm:rz}
Let $\K$ be any field of characteristic $\neq 2$, and for any $r, s \ge 1$ let $E_{(a,b),(c,d)} \in M_{rs}(\K)$ denote the standard matrix unit with $1$ at entry $((a,b),(c,d))$ and $0$ elsewhere.
\begin{enumerate}
\item \textbf{Refutation of Lemma~3.13} (\lean{rz\_lemma313\_counterexample}, \lean{rzM1\_hasSoS\_two}): Let $M_1 := E_{(0,0),(0,1)} \in M_4(\K)$ (indexed by $\mathrm{Fin}\,2 \times \mathrm{Fin}\,2$), so $Q_{M_1}(x,y) = x_0^2 y_0 y_1$. Then $\PT_\K(M_1) = 1$, whereas $\SoS_\K(Q_{M_1}) > 1$. When $\sqrt{-1} \in \K$, $\SoS_\K(Q_{M_1}) = 2$.
\item \textbf{Refutation of the left inequality of Proposition~3.16} (\lean{rz\_prop316\_left\_counterexample}, \lean{rzM2\_hasSoS\_four}): Let $M_2 := E_{(0,0),(0,1)} + E_{(1,2),(1,3)} \in M_{16}(\K)$ (indexed by $\mathrm{Fin}\,4 \times \mathrm{Fin}\,4$), so $Q_{M_2}(x,y) = x_0^2 y_0 y_1 + x_1^2 y_2 y_3$. Then $\PT_\K(M_2) \le 2$, whereas $\SoS_\K(Q_{M_2}) > 3$ (so $\SoS_\K(Q_{M_2}) - 1 \ge 3 > 2 \ge \PT_\K(M_2)$). When $\sqrt{-1} \in \K$, $\SoS_\K(Q_{M_2}) = 4$.
\item \textbf{Refutation of the right inequality of Proposition~3.16} (\lean{rz\_prop316\_right\_counterexample}): Let $M_3 := E_{(0,0),(0,1)} - E_{(0,1),(0,0)} \in M_4(\K)$. Then $Q_{M_3} \equiv 0$, so $\SoS_\K(Q_{M_3}) = 0$, whereas $M_3 \neq 0$ implies $\PT_\K(M_3) \ge 1 > 2 \cdot 0$.
\end{enumerate}
In particular, the repaired inequality $\SoS_\K(Q_M) \le 2\,\PT_\K(M)$ of Theorem~\ref{thm:dict}(5) is tight for general $M$ (realised by $M_1$ and $M_2$, and by direct sums of $k$ copies of $M_1$ with $\PT = k$ and $\SoS = 2k$).
\end{theorem}

\begin{proof}
(1) $M_1$ has rank $1$, so $\PT_\K(M_1) \le 1$. If $Q_{M_1} = L^2$ for a single bilinear form $L$, comparing the $(0,0),(0,0)$ entry of $\sym$ gives $4 L_{(0,0)}^2 = 0$, hence $L_{(0,0)} = 0$; then comparing $(0,0),(0,1)$ gives $0 = 4 L_{(0,0)} L_{(0,1)} = 2$, contradicting $\operatorname{char}\K \neq 2$. With $i^2 = -1$, $x_0^2 y_0 y_1 = (\frac{1}{2} x_0(y_0+y_1))^2 + (\frac{i}{2} x_0(y_0-y_1))^2$ gives $\SoS_\K(Q_{M_1}) = 2$.

(2) $M_2$ is a sum of two rank-$1$ matrices, so $\PT_\K(M_2) \le 2$. Suppose $Q_{M_2} = \sum_{k=1}^3 L_k^2$, and write $c_{(a,b)} \in \K^3$ for the column $(L_k(a,b))_{k=1}^3$. Extracting entries of $\sym(\sum \ell_k \ell_k^\top) = \sym M_2$ gives $c_{(a,b)} \cdot c_{(a,b)} = 0$ for all $(a,b)$, and $4\,c_{(a,b)} \cdot c_{(a,d)} = (\sym M_2)_{(a,b),(a,d)}$. Hence $u_1 := c_{(0,0)}, u_2 := c_{(0,1)}$ form a hyperbolic plane ($u_1^2 = u_2^2 = 0$, $u_1\cdot u_2 \neq 0$), and $v_1 := c_{(1,2)}, v_2 := c_{(1,3)}$ form a second hyperbolic plane ($v_1^2 = v_2^2 = 0$, $v_1\cdot v_2 \neq 0$). Moreover, $c_{(0,2)}$ and $c_{(0,3)}$ are isotropic and orthogonal to both $u_1$ and $u_2$. In $\K^3$ with the standard nondegenerate dot product, any isotropic vector $w$ orthogonal to a hyperbolic plane $\{u_1, u_2\}$ must vanish: otherwise $\{u_1, u_2, w\}$ would be linearly independent, hence a basis of $\K^3$, forcing $w$ orthogonal to all of $\K^3$ (private helper \lean{zero\_of\_iso\_orth\_hyperbolic}). Thus $c_{(0,2)} = c_{(0,3)} = 0$. The cross-entries $(0,b),(1,d)$ for $b \in \{0,1\}$, $d \in \{2,3\}$ give $2(c_{(0,b)}\cdot c_{(1,d)} + c_{(0,d)}\cdot c_{(1,b)}) = 0$, which reduces to $u_i \cdot v_j = 0$ for all $i,j \in \{1,2\}$. Hence $\{u_1, u_2, v_1, v_2\}$ are four mutually orthogonal hyperbolic vectors, therefore linearly independent in $\K^3$, contradiction (private helper \lean{no\_two\_orth\_hyperbolic\_in\_fin3}). With $i^2 = -1$, polarising both terms yields $\SoS_\K(Q_{M_2}) = 4$.

(3) Direct entrywise check gives $\sym M_3 = 0$, so $Q_{M_3} = 0$ has a $0$-square representation, while $(M_3)_{(0,0),(0,1)} = 1 \neq 0$ implies $\PT_\K(M_3) \ge 1$.
\end{proof}

\begin{remark}[Full symmetry and the open question]
The proof of Lemma~3.13 in~\cite{RZ25} sets $D = 2^{-d}\sum_{\kappa} P_\kappa$, which is not symmetric in general even when $M$ is \emph{fully symmetric}. Note that $M_1, M_2, M_3$ are not fully symmetric (for $M_1$, $(M_1)_{(0,0),(0,1)} = 1 \neq 0 = (M_1)_{(0,1),(0,0)}$). For fully symmetric $M$, whether $\SoS_\C(Q_M) \le \PT_\C(M)$ holds---and in particular whether $\SoS_\C(X_{r,s}) \le \TC_\C(r,s)$---remains open. By Theorem~\ref{thm:dict}(5), $\SoS_\C(X_{r,s}) \le 2\TC_\C(r,s)$ holds unconditionally.
\end{remark}

\section{Characteristic-free crossing bounds}\label{sec:lb}

We now determine fundamental lower and upper bounds on $\TC_\K(r,s) = \TC_\K(I_{rs})$ that hold over \emph{every} field $\K$, including $\operatorname{char}\K = 2$.

\begin{theorem}[Crossing bounds and exact small values; \lean{LowerBounds.lean}]\label{thm:lb}
Let $\K$ be any field (arbitrary characteristic).
\begin{enumerate}
\item \textbf{Slice lower bound} (\lean{twoCut\_ge\_max}): for all $r, s \ge 1$, $\TC_\K(r,s) \ge \max(r, s)$.
\item \textbf{Strict beating of $\max(r,s)$ for $r, s \ge 3$} (\lean{twoCut\_gt\_max\_of\_ge\_three}): for all $r, s \ge 3$, $\TC_\K(r,s) \ge \max(r, s) + 1$.
\item \textbf{Exact $2 \times s$ formula} (\lean{twoCut\_two\_iff}): for all $s \ge 1$, $\TC_\K(2, s) = 2\lceil s/2\rceil$.
\item \textbf{Transpose upper bound} (\lean{transpose\_twoCut}): for all $r, s \ge 1$ with $m = \min(r,s)$,
\[
  \TC_\K(r,s) \le 1 + rs - \frac{m(m+1)}{2}.
\]
\item \textbf{Exact $3 \times 3$ value} (\lean{twoCut\_three\_three\_iff}): $\TC_\K(3, 3) = 4$.
\end{enumerate}
\end{theorem}

\begin{proof}[Proof of (1) and (2)]
Let $I_{rs} = A + W^\Gamma$ with $a = \rk A$, $b = \rk W$, and suppose without loss of generality $r \le s$.

For (1), fix any row-slice $x \in \K^r$ and column-slice $\varphi \in \K^r$ with $\varphi^\top x = 1$, and form the $s\times s$ contracted matrices $A_{x,\varphi} := (x^\top \otimes I_s) A (\varphi \otimes I_s)$ and $(W^\Gamma)_{x,\varphi}$. Contracting $I_{rs} = A + W^\Gamma$ yields
\[
  I_s = A_{x,\varphi} + (W^\Gamma)_{x,\varphi}, \qquad
  \rk A_{x,\varphi} \le a, \quad \rk(W^\Gamma)_{x,\varphi} \le b.
\]
Indeed, for $W = \sum_{t=1}^b p_t q_t^\top$, one computes
\[
  ((W^\Gamma)_{x,\varphi})_{\alpha,\beta}
  = \sum_{i,k} x_i \varphi_k W_{(i,\beta),(k,\alpha)}
  = \sum_{t=1}^b \Bigl(\sum_i x_i (p_t)_{(i,\beta)}\Bigr)\Bigl(\sum_k \varphi_k (q_t)_{(k,\alpha)}\Bigr),
\]
which is a sum of $b$ rank-$1$ matrices. Hence $s = \rk I_s \le a + b$. By symmetry of $r$ and $s$, $\max(r,s) \le a + b$.

For (2), suppose for contradiction that $r, s \ge 3$ and $a + b = s$ (with $r \le s$). Since $I_s = A_{x,\varphi} + (W^\Gamma)_{x,\varphi}$ with $\rk A_{x,\varphi} + \rk(W^\Gamma)_{x,\varphi} \le a + b = s$, every slice $(x,\varphi)$ with $\varphi^\top x = 1$ must attain equality $\rk A_{x,\varphi} = a$ and $\rk(W^\Gamma)_{x,\varphi} = b$, and moreover $P_{x,\varphi} := A_{x,\varphi}$ is an idempotent ($P_{x,\varphi}^2 = P_{x,\varphi}$). Writing $A = U V^\top$ with $U, V \in \K^{(r\times s)\times a}$ and splitting $U, V$ into $r$ blocks $U_i, V_i \in \K^{s\times a}$ ($i \in [r]$), we have $A_{x,\varphi} = U_x V_\varphi^\top$ where $U_x = \sum_i x_i U_i$ and $V_\varphi = \sum_i \varphi_i V_i$. Full rank $a$ of $U_x V_\varphi^\top$ together with $(U_x V_\varphi^\top)^2 = U_x V_\varphi^\top$ forces $V_\varphi^\top U_x = I_a$ for every $(x,\varphi)$ with $\varphi^\top x = 1$. Specialising to $(e_i, e_i)$ and $(e_i, e_i + e_k)$ yields the block biorthogonality
\[
  V_k^\top U_i = \delta_{k,i}\, I_a \qquad (i, k \in [r]).
\]
This has two characteristic-free consequences:
\begin{enumerate}
\item[(i)] Stacking the blocks horizontally gives $V_{\mathrm{all}} U_{\mathrm{all}} = I_{ra}$, so $ra \le s$.
\item[(ii)] Define the linear map $\Psi : \ker(W) \to \K^{\{(k,i) : k \le i\} \times [a]}$ by
\[
  (\Psi v)_{(k,i),t} := (U_k^\top v_i)_t = \sum_{\beta \in [s]} (U_k)_{\beta, t}\, v_{(i,\beta)}.
\]
We claim $\Psi$ is injective. For $v \in \ker(W)$, $W = I - A^\Gamma$ gives $v = A^\Gamma v$, i.e.\ in block coordinates
\[
  v_{(k,\beta)} = \sum_{i \in [r]} \sum_{t' \in [a]} (V_i)_{\beta, t'}\, M(k, i, t'), \qquad
  M(k, i, t) := \sum_{\alpha \in [s]} (U_k)_{\alpha, t}\, v_{(i,\alpha)}.
\]
Substituting this expression into $M(j, k, t)$ and using $V_i^\top U_j = \delta_{i,j} I_a$ immediately yields $M(j, k, t) = M(k, j, t)$ for all $j, k, t$. Notice that no division by $2$ and no trace is used anywhere. Thus $\Psi(v) = 0$ implies $M(k, i, t) = 0$ for all $k \le i$, hence by symmetry for all $k, i$, hence $v = 0$.
\end{enumerate}
By injectivity of $\Psi$, $\dim\ker(W) \le \binom{r+1}{2} a$. Since $\dim\ker(W) = rs - b$ and $b = s - a$, we obtain
\[
  2(r-1)s \le (r^2 + r - 2) a.
\]
Combined with $ra \le s$ (hence $2(r-1)ra \le (r^2 + r - 2)a$), if $a > 0$ then $2r(r-1) \le r^2 + r - 2$, i.e.\ $(r-1)(r-2) \le 0$, contradicting $r \ge 3$. And if $a = 0$, then $b = s$ gives $W = I^\Gamma = I$, whose rank is $rs > s$ (since $r \ge 3, s \ge 1$), contradiction.
\end{proof}

\begin{proof}[Proof of (3)--(5)]
For (3), the upper bound $\TC_\K(2, s) \le 2\lceil s/2\rceil$ is achieved by taking $\lceil s/2\rceil$ disjoint $2\times 2$ blocks, each decomposed by the $(2,2)$ transpose decomposition $I_4 = (I_4 - \mathrm{SWAP}) + \mathrm{SWAP}$ of cost $1 + 1 = 2$ (truncating the last column when $s$ is odd). For the lower bound when $s$ is odd, if $\TC_\K(2, s) = s$ with $s = 2k+1$, the same block biorthogonality analysis at $r = 2$ gives $2a \le s$ and $2s \le 4a$ (from $2(r-1)s \le (r^2+r-2)a$ at $r=2$), forcing $2a = s$, impossible for odd $s$.

For (4), let $m = \min(r,s)$ and set $W = u u^\top$ where $u_{(a,b)} = \mathbf{1}_{a = b < m}$. Then $\rk W = 1$, and $W^\Gamma$ fixes $(i,i)$ ($i < m$) and swaps $(i,j) \leftrightarrow (j,i)$ ($i < j < m$). The $m(m+1)/2$ vectors $e_{(i,i)}$ and $e_{(i,j)} + e_{(j,i)}$ ($i < j < m$) are linearly independent in $\ker(I - W^\Gamma)$ over \emph{every} field (including $\operatorname{char}\K = 2$), so $\rk(I - W^\Gamma) \le rs - m(m+1)/2$.

For (5), combining (2) ($4 \le \TC_\K(3,3)$) and (4) ($\TC_\K(3,3) \le 1 + 9 - 6 = 4$) gives $\TC_\K(3,3) = 4$.
\end{proof}

\section{Explicit exact and border constructions}\label{sec:witnesses}

\begin{theorem}[Constructions and border decompositions; \lean{Witnesses.lean}]\label{thm:witnesses}
Let $\K$ be a field.
\begin{enumerate}
\item \textbf{Clifford self-dual bound} (\lean{clifford\_twoCut}): If $\operatorname{char}\K \neq 2$ and $\gamma_1,\dots,\gamma_n \in M_d(\K)$ satisfy $\gamma_i \gamma_j + \gamma_j \gamma_i = 2\delta_{ij} I_d$, then $A_{(a,i),(c,j)} := \frac{1}{2}(\gamma_i \gamma_j)_{a,c}$ satisfies $A + A^\Gamma = I_{dn}$ and $\rk A \le d$, so $\TC_\K(d,n) \le 2d$.
\item \textbf{Border-to-exact conversion} (\lean{exact\_of\_border\_coeffs}): Over \emph{every} field $\K$, if a pole-order-$k$ formal Laurent family $P(\varepsilon) = \varepsilon^{-k} U(\varepsilon) V(\varepsilon)^\top$ (width $a$) and $Q(\varepsilon) = \varepsilon^{-k} X(\varepsilon) Y(\varepsilon)^\top$ (width $b$) has $\varepsilon^0$-coefficient $\sum_{i=0}^k U_i V_{k-i}^\top + \bigl(\sum_{i=0}^k X_i Y_{k-i}^\top\bigr)^\Gamma = I_{rs}$, then $\TC_\K(r,s) \le (k+1)(a+b)$.
\item \textbf{Characteristic-free $\pm 1$ exact schemes} (\lean{tc45\_twoCut}, \lean{tc46\_twoCut}, \lean{tc88\_twoCut}): Over \emph{every} field $\K$ (including $\operatorname{char}\K = 2$),
\[
  \TC_\K(4,5) \le 8, \qquad \TC_\K(4,6) \le 8, \qquad \TC_\K(8,8) \le 16.
\]
\item \textbf{Pole-order-$2$ border decomposition of $(4,4)$} (\lean{border44\_rank\_le}, \lean{border44\_laurent\_identity}, \lean{border44\_identity\_field}): Over \emph{every} field $\K$, $I_{16}$ admits an explicit integer pole-order-$2$ border two-cut decomposition of ranks $3 + 3$ ($\bTC_{\K,2}(4,4) \le 6$): the explicit width-$3$ polynomial matrices $P(\varepsilon), Q(\varepsilon)$ satisfy $\rk P(\varepsilon) \le 3$, $\rk Q(\varepsilon) \le 3$, and $P(\varepsilon) + Q(\varepsilon)^\Gamma = \varepsilon^2 I_{16} + \varepsilon^3 R_1 + \varepsilon^4 R_2$ for every $\varepsilon \in \K$ (including $\K(t)$).
\end{enumerate}
\end{theorem}

\begin{remark}[Small-size table]
Table~\ref{tab:small} summarises the exact and border bounds on $\TC_\K(r,s)$ proved in this paper across all characteristics.
\end{remark}

\begin{table}[htbp]
\centering
\small
\begin{tabular}{@{}lccc>{\raggedright\arraybackslash}p{5.8cm}@{}}
\toprule
$(r, s)$ & Lower & Upper & $\bTC_2$ & Lean witnesses \\
\midrule
$(2, s)$ & $2\lceil s/2\rceil$ & $2\lceil s/2\rceil$ & --- & \lean{twoCut\_two\_iff} \\
$(3, 3)$ & $4$ & $4$ & --- & \lean{twoCut\_three\_three\_iff} \\
$(4, 4)$ & $5$ & $7$ & $\le 6$ & \lean{twoCut\_gt\_max\_of\_ge\_three}, \lean{transpose\_twoCut}, \lean{border44\_rank\_le}, \lean{border44\_laurent\_identity} \\
$(4, 5)$ & $6$ & $8$ & --- & \lean{twoCut\_gt\_max\_of\_ge\_three}, \lean{tc45\_twoCut} \\
$(4, 6)$ & $7$ & $8$ & --- & \lean{twoCut\_gt\_max\_of\_ge\_three}, \lean{tc46\_twoCut} \\
$(8, 8)$ & $9$ & $16$ & --- & \lean{twoCut\_gt\_max\_of\_ge\_three}, \lean{tc88\_twoCut} \\
\bottomrule
\end{tabular}
\caption{Characteristic-free bounds on $\TC_\K(r,s)$ proved in \lean{PastaMath.TwoCutSoS}.}
\label{tab:small}
\end{table}

\section{Reflection-type word-span tensor powers}\label{sec:refl}

A recurring strategy for improving asymptotic exponents from a base decomposition $I_D = M + N$ ($D = rs$, $N = I_D - M$) is to take $K$-th tensor powers and form a \emph{word-span decomposition}
\[
  P = F(M_1, \dots, M_K) = \sum_{w \in \{0,1\}^K} c_w \prod_{i=1}^K M_i^{w_i} N_i^{1-w_i},
\]
where $M_i = I^{\otimes (i-1)} \otimes M \otimes I^{\otimes (K-i)}$ and $N_i = I - M_i$. A $0/1$ choice $c_w \in \{0,1\}$ is a \emph{grouping}; arbitrary scalars $c_w \in \R$ allow algebraic cancellation.

\subsection{The $p$-biased reflection operator $T_{-1}$}

\begin{definition}[Reflection-type spectrum]
A matrix $M \in M_D(\R)$ ($D = rs$) has \emph{reflection-type spectrum} with multiplicity $0 < n < D$ and eigenvalue $\lambda$ if $M$ is similar to $\operatorname{diag}(0^{D-n}, \lambda^n)$ and $M^\Gamma$ is similar to $\operatorname{diag}(1^{D-n}, (1-\lambda)^n)$.
\end{definition}

Many canonical exact decompositions have reflection-type spectra: Hurwitz $[n,n,n]$ for $n \in \{2,4,8\}$ ($p = 1/n$), Clifford self-dual decompositions with $s$ generators ($p = 1/s$), the $(2,2)$ transpose point ($p = 1/4$), and our $(4,6)$ scheme ($p = 1/6$).

\begin{proposition}[Trace identity, slice constraint, and $T_{-1}$; \lean{ReflectionPowers.lean}]\label{prop:refl}
Suppose $M \in M_{rs}(\R)$ has reflection-type spectrum with parameters $(n, \lambda)$, and put $p := n/D \in (0,1)$.
\begin{enumerate}
\item (\lean{reflection\_lam\_eq}) Since $\tr M = \tr(M^\Gamma)$, necessarily $2\lambda n = D$, i.e.\ $\lambda = 1/(2p)$.
\item (\lean{reflection\_slice\_constraint}) Both $M$ and $\frac{\lambda}{2\lambda-1}(I - \lambda^{-1} M)$ give two-cut decompositions of $I_{rs}$ (of costs $\le 2n$ and $\le 2(D-n)$), so Theorem~\ref{thm:lb}(1) forces the \emph{slice constraint}
\[
  \max(r,s) \le 2\min(n, D - n), \qquad\text{equivalently}\qquad \min(p, 1-p) \ge \frac{1}{2\min(r,s)}.
\]
\item (\lean{reflOp\_involutive}, \lean{reflOp\_selfadjoint}, \lean{reflOp\_const\_one}, \lean{wordPoly\_eval\_B}) Let $\mu_p(x) := \prod_{i=1}^K p^{x_i}(1-p)^{1-x_i}$ on $\{0,1\}^K$, and define $E_1(p) := \bigl(\begin{smallmatrix} 1-2p & 2p \\ 2-2p & 2p-1 \end{smallmatrix}\bigr)$ and $T_{-1} := E_1(p)^{\otimes K}$. Then $T_{-1}$ is an involution, is self-adjoint in $L^2(\mu_p)$, and fixes $\mathbf{1}$. For every multilinear word polynomial $F$ with $A$-eigenvalue function $h(a) := F(\lambda a_1, \dots, \lambda a_K)$, the $B$-eigenvalue function satisfies $F(1 - \lambda b_1, \dots, 1 - \lambda b_K) = (T_{-1} h)(b)$. Consequently, in the simultaneous eigenbases of $M^{\otimes K}$ and $(M^\Gamma)^{\otimes K}$, the normalised word-span cost of $h$ is
\[
  \varphi_K(p; h) := \mu_p(\supp h) + \mu_p(\supp(\mathbf{1} - T_{-1} h)),
\]
and $\varphi_K(1-p; h) = \varphi_K(p; h \circ \operatorname{flip})$ (\lean{wordSpanCost\_one\_sub}).
\end{enumerate}
\end{proposition}

\subsection{The Valuation Theorem: when groupings cannot cancel}

Without algebraic cancellation, each word $w \in \{0,1\}^K$ is assigned either to the $A$-side or the $B$-side, paying at least the Bayes error $\sum_{T \in \{0,1\}^K} \min(\mu_p(T), \mu_{1-p}(T))$, whose $K$-th root tends to the Bhattacharyya coefficient $2\sqrt{p(1-p)}$. When $\lambda = 2$ ($p = 1/4$), cancellation beats this threshold (for instance at $K = 4$, one can achieve $62/256$ versus the Bayes error $80/256 = 5/16$). Our Valuation Theorem explains exactly when no integral or $\ell$-local cancellation can occur.

\begin{theorem}[Valuation Theorem; \lean{valuation\_theorem}, \lean{layer\_grouping\_le}]\label{thm:val}
Let $\lambda = a/b \in \Q$ ($a,b \in \Z \setminus \{0\}$) with $p = b/(2a) \in [0,1]$, and suppose there is a prime $\ell \mid (a-b)$ with $\ell \nmid a$ (in particular, this holds whenever $a/b$ is in lowest terms and $|a-b| \ge 2$). Then for every $K \ge 1$ and every word-span polynomial $F$ with $\ell$-integral rational coefficients $c_w \in \Z_{(\ell)}$ (including every $0/1$ grouping and every integer combination):
\[
  \varphi_K(p; h) \ge \sum_{T \in \{0,1\}^K} \min\bigl(\mu_p(T),\, \mu_{1-p}(T)\bigr).
\]
Moreover, for $0 < p \le 1/2$, the $0/1$ layer grouping $c_w = \mathbf{1}_{2|w| \ge K}$ attains equality at every $K \ge 1$.
\end{theorem}

\begin{proof}
Since $\ell \mid (a-b)$ and $\ell \nmid a$, also $\ell \nmid b$. Let $D_c = \prod_w \operatorname{den}(c_w)$, so $\ell \nmid D_c$, and write $c_w = \tilde c_w / D_c$ with $\tilde c_w \in \Z$. Fix $T \in \{0,1\}^K$. Clearing denominators by $D_c b^K$:
\begin{itemize}
\item $D_c b^K h(T) = \sum_{w} \tilde c_w \prod_{i=1}^K \bigl(w_i \cdot (a T_i) + (1-w_i)(b - a T_i)\bigr)$. For every $w \neq T$, some coordinate $i$ has either $(w_i, T_i) = (1, 0)$ (contributing $0$) or $(w_i, T_i) = (0, 1)$ (contributing $b - a \equiv 0 \pmod{\ell}$). Hence $D_c b^K h(T) \equiv \tilde c_T a^{|T|} b^{K-|T|} \pmod{\ell}$.
\item Writing $1 - (T_{-1} h)(\neg T)$ in the same basis gives $D_c b^K (1 - (T_{-1} h)(\neg T)) \equiv (D_c - \tilde c_T) a^{K-|T|} b^{|T|} \pmod{\ell}$.
\end{itemize}
Since $\ell \nmid a$ and $\ell \nmid b$, if both $h(T) = 0$ and $(T_{-1} h)(\neg T) = 1$, then $\ell \mid \tilde c_T$ and $\ell \mid (D_c - \tilde c_T)$, so $\ell \mid D_c$, contradiction. Pairing $T$ on the $A$-side with $\neg T$ on the $B$-side (which has weight $\mu_p(\neg T) = \mu_{1-p}(T)$) yields $\varphi_K(p; h) \ge \sum_T \min(\mu_p(T), \mu_{1-p}(T))$.

For attainment when $0 < p \le 1/2$, take $c_w = \mathbf{1}_{2|w| \ge K}$. If $2|T| < K$, every $w \le T$ has $2|w| < K$, so $h(T) = 0$ and the pair $(T, \neg T)$ pays at most $\mu_{1-p}(T) = \min(\mu_p(T), \mu_{1-p}(T))$. If $2|T| \ge K$, every $w \ge T$ has $2|w| \ge K$, so $1 - c_w = 0$, hence $(T_{-1} h)(\neg T) = 1$, and the pair pays at most $\mu_p(T) = \min(\mu_p(T), \mu_{1-p}(T))$.
\end{proof}

Notice that Theorem~\ref{thm:val} immediately covers Hurwitz $[8,8,8]$ ($\lambda = 4$, $\ell = 3$), our $(4,6)$ scheme ($\lambda = 3$, $\ell = 2$), Clifford $(4,5)$ ($\lambda = 5/2$, $\ell = 3$), and Clifford with any $s \ge 5$ generators ($\lambda = s/2$).

\subsection{Tensorising certificate packings and the no-go theorem}

To lower-bound $\varphi_K(p; h)$ for \emph{arbitrary} real coefficients $c_w \in \R$ (even when cancellation occurs, notably at $p = 1/4$), we introduce a dual packing framework that tensorises cleanly across block lengths.

\begin{definition}[Farkas certificates and packings; \lean{ReflCert}, \lean{CertPacking}]
A \emph{certificate} at block length $K$ for weight $p$ is a vector $y : \{0,1\}^K \to \R$ with $\sum_b y(b) \neq 0$; its $A$-grid image is $y_A := (E_1(p)^\top)^{\otimes K} y$. A \emph{certificate packing} $P = \{(y^{(i)}, w_i)\}_{i \in I}$ assigns weights $w_i \ge 0$ such that
\[
  \sum_{i : y_A^{(i)}(a) \neq 0} w_i \le \mu_p(a) \quad (\forall a), \qquad
  \sum_{i : y^{(i)}(b) \neq 0} w_i \le \mu_p(b) \quad (\forall b).
\]
Its value is $\operatorname{val}(P) := \sum_{i \in I} w_i$.
\end{definition}

\begin{lemma}[Packing bound and tensorisation; \lean{CertPacking}]\label{lem:pack}
\begin{enumerate}
\item For every certificate $y$ and every $h : \{0,1\}^K \to \R$, either $\supp(y_A) \cap \supp(h) \neq \emptyset$ or $\supp(y) \cap \supp(\mathbf{1} - T_{-1} h) \neq \emptyset$.
\item Consequently, $\operatorname{val}(P) \le \varphi_K(p; h)$ for every packing $P$ and every $h$.
\item If $P_1, P_2$ are packings at block lengths $K_1, K_2$, their tensor product $P_1 \otimes P_2$ (with $y^{(i,j)} = y^{(i)} \otimes y^{(j)}$ and $w_{i,j} = w_i w_j$) is a packing at $K_1 + K_2$ of value $\operatorname{val}(P_1)\operatorname{val}(P_2)$.
\item Hence any packing $P$ at block length $B \ge 1$ with $\operatorname{val}(P) > 0$, together with a positive $K=1$ packing, yields a constant $c > 0$ such that $c \cdot \operatorname{val}(P)^{K/B} \le \varphi_K(p; h)$ for all $K \ge 1$ and all $h$.
\end{enumerate}
\end{lemma}

We construct four explicit packings in \lean{ReflectionPowers.lean}:
\begin{itemize}
\item \textbf{Trivial $K=1$ packing} (\lean{packing\_K1\_trivial}): single certificate $\delta_0$ with weight $\min(p, 1-p) > 0$ for all $p \in (0,1)$.
\item \textbf{Parametric $K=1$ packing} (\lean{packing\_K1\_mid}): four certificates $a = \delta_0$, $b = \delta_1$, $c = (2-2p, 2p-1)$, $d = (1-2p, 2p)$ with weights $(2/3-p, p-1/3, p-1/3, 2/3-p)$ on $p \in [1/3, 2/3]$, attaining value $2/3$.
\item \textbf{Parametric $K=2$ packing} (\lean{packing\_K2}): three certificates $d\otimes a$, $a\otimes d$, and the eigenvector $C_3 = (1-p)\delta_{00} + p\delta_{11}$ (which satisfies $(E_1^\top \otimes E_1^\top) C_3 = C_3$), with weights $\alpha(p(1-p), p(1-p), p^2)$ where $\alpha = \min(1, (1-p)^2/(2p-p^2))$, attaining value $\min(2p-p^2, (1-p)^2)$ on $p \in (0, 1/2]$.
\item \textbf{Exact $22$-certificate $K=4$, $p=1/4$ packing} (\lean{packing\_K4\_quarter}): $22$ rational certificates (scaled by $24$ to integers in $[-36, 18]$) with rational weights (scaled by $2048$ to positive integers summing to $456$), kernel-checked by \lean{decide} to satisfy all $32$ load constraints and attain exact value $456/2048 = 57/256$.
\end{itemize}

\begin{theorem}[Reflection-type word-span no-go theorem; \lean{reflection\_wordSpan\_nogo}]\label{thm:nogo}
Let $r, s \ge 2$, $D = rs$, and $0 < n < D$ satisfy the slice constraint $\max(r,s) \le 2\min(n, D-n)$ (which holds automatically for every exact reflection-type matrix decomposition over $\R$ by Proposition~\ref{prop:refl}(2)). Put $p = n/D$. Then there exist $c > 0$ and $g > 0$ such that
\[
  \log(D g) > \sqrt{(\log r)^2 + (\log s)^2}
\]
and, for every tensor power $K \ge 1$ and every word-span evaluation $h : \{0,1\}^K \to \R$,
\[
  c \cdot (D g)^K \le D^K \varphi_K(p; h).
\]
In particular, no reflection-type decomposition yields word-span tensor powers with asymptotic rate $\le \sqrt{1 + q^2}$ where $q = \log\min(r,s)/\log\max(r,s)$.
\end{theorem}

\begin{proof}[Proof outline]
By $\varphi_K(1-p; h) = \varphi_K(p; h \circ \operatorname{flip})$, assume $p \le 1/2$, so $1/(2m) \le p \le 1/2$ where $m = \min(r,s)$ and $M = \max(r,s)$.
\begin{itemize}
\item \textbf{Case $(r,s) = (2,2)$.} Then $n \in \{1, 2\}$, so $p \in \{1/4, 1/2\}$. At $p = 1/4$, \lean{packing\_K4\_quarter} gives $g = (57/256)^{1/4}$, and $\log(4g) = \frac{1}{4}\log 57 > \sqrt{2}\log 2$ since $2^{4\sqrt{2}} < 2^{17/3} < 57$ (\lean{k4\_beats}, via $2^{17} = 131072 < 185193 = 57^3$). At $p = 1/2$, \lean{packing\_K1\_mid} gives $g = 2/3 > 2^{\sqrt{2}-2}$ (\lean{two\_thirds\_gt}, via $\sqrt{2} < 58/41$).
\item \textbf{Case $m = 2$, $M \ge 3$.} Here $p \ge 1/4$. Combining \lean{packing\_K2} on $[1/4, 1/3]$ (value $\ge 7/16$) and \lean{packing\_K1\_mid} on $[1/3, 1/2]$ (value $2/3 > \sqrt{7}/4$) gives $g \ge \sqrt{7}/4$. Writing $L = \log M \ge \log 3 > 1$, $\log(7/4) > 1/2$, and $\log 2 < 7/10$, we get $(L + \frac{1}{2}\log(7/4))^2 > L^2 + (\log 2)^2$, hence $\log(2Mg) > \sqrt{L^2 + (\log 2)^2}$.
\item \textbf{Case $m \ge 3$.} Here $p \ge 1/(2m)$. Combining \lean{packing\_K2} on $[1/(2m), 1/3]$ and \lean{packing\_K1\_mid} on $[1/3, 1/2]$ gives $g \ge \sqrt{(4m-1)/(4m^2)}$ (since $(4m-1)/(4m^2) \le 4/9$ for $m \ge 3$). Writing $L = \log M \ge \ell = \log m > 0$, the elementary bound $\sqrt{L^2+\ell^2} \le L + (\sqrt{2}-1)\ell$ reduces the desired inequality to $m^{\sqrt{2}-2} < \sqrt{(4m-1)/(4m^2)}$, which follows from $\sqrt{2} < 17/12$ and $4096 m^5 < (4m-1)^6$ for $m \ge 3$ (\lean{selfdual\_rate\_m3}). \qedhere
\end{itemize}
\end{proof}

\section{Formalisation and verification}\label{sec:verified}

The Lean~4 formalisation resides in \lean{PastaMath/TwoCutSoS/} (Lean~4.31.0, Mathlib revision \lean{fabf563a7c95a166b8d7b6efca11c8b4dc9d911f}). Every headline theorem is audited by \lean{PastaMath/TwoCutSoS/Axioms.lean} and depends exclusively on the standard axioms \lean{[propext, Classical.choice, Quot.sound]} with zero \lean{sorry}, zero \lean{native\_decide}, and zero custom axioms.

\begin{table}[htbp]
\centering
\small
\begin{tabular}{@{}ll>{\raggedright\arraybackslash}p{6.6cm}@{}}
\toprule
Paper result & Lean module & Headline Lean theorems \\
\midrule
Theorem~\ref{thm:dict} (Dictionary) & \lean{Defs.lean}, \lean{Dictionary.lean} & \lean{sameBiquad\_iff\_qf}, \lean{sameBiquad\_iff\_sym4}, \lean{hasPTRank\_iff\_hasTwoCut}, \lean{prod\_of\_twoCut}, \lean{twoCut\_of\_sos}, \lean{twoCut\_of\_prod}, \lean{sos\_of\_prod}, \lean{sos\_of\_twoCut}, \lean{sos\_of\_ptRank} \\
Theorem~\ref{thm:rz} (RZ \S3.2) & \lean{RossmanZhu.lean} & \lean{rz\_lemma313\_counterexample}, \lean{rzM1\_hasSoS\_two}, \lean{rz\_prop316\_left\_counterexample}, \lean{rzM2\_hasSoS\_four}, \lean{rz\_prop316\_right\_counterexample} \\
Theorem~\ref{thm:lb} (Crossing) & \lean{LowerBounds.lean} & \lean{twoCut\_ge\_max}, \lean{twoCut\_gt\_max\_of\_ge\_three}, \lean{twoCut\_two\_iff}, \lean{transpose\_twoCut}, \lean{twoCut\_three\_three\_iff} \\
Theorem~\ref{thm:witnesses} (Witnesses) & \lean{Witnesses.lean} & \lean{clifford\_twoCut}, \lean{exact\_of\_border\_coeffs}, \lean{tc45\_twoCut}, \lean{tc46\_twoCut}, \lean{tc88\_twoCut}, \lean{border44\_rank\_le}, \lean{border44\_laurent\_identity}, \lean{border44\_identity\_field} \\
Proposition~\ref{prop:refl} ($T_{-1}$) & \lean{ReflectionPowers.lean} & \lean{reflOp\_involutive}, \lean{reflOp\_selfadjoint}, \lean{wordPoly\_eval\_B}, \lean{reflection\_slice\_constraint} \\
Theorem~\ref{thm:val} (Valuation) & \lean{ReflectionPowers.lean} & \lean{valuation\_theorem}, \lean{layer\_grouping\_le} \\
Lemma~\ref{lem:pack}, Thm.~\ref{thm:nogo} & \lean{ReflectionPowers.lean} & \lean{CertPacking.lower\_bound}, \lean{CertPacking.asymptotic\_bound}, \lean{packing\_K4\_quarter\_value}, \lean{reflection\_wordSpan\_nogo} \\
\bottomrule
\end{tabular}
\caption{Mapping between paper statements and kernel-checked Lean~4 theorems.}
\label{tab:lean}
\end{table}

\paragraph{What is not machine-checked.}
\begin{enumerate}
\item The Kronecker-product simultaneous-diagonalisation step identifying $\rk F(M_1,\dots,M_K) + \rk(I - F(M_1,\dots,M_K))^\Gamma = D^K \varphi_K(p; h)$ for a diagonalisable reflection-type base $M$ is proved on paper; Lean kernel-checks the slice constraint (\lean{reflection\_slice\_constraint}), the eigenvalue transformation (\lean{wordPoly\_eval\_B}), the Valuation Theorem (\lean{valuation\_theorem}, \lean{layer\_grouping\_le}), the certificate packings, and the asymptotic lower bound on $D^K \varphi_K(p; h)$ (\lean{reflection\_wordSpan\_nogo}).
\item Theorem~\ref{thm:nogo} applies to word-span decompositions of reflection-type diagonalisable bases over $\R$; it does not cover decompositions outside the word algebra, non-reflection spectra, non-diagonalisable bases, or border families.
\item As noted in Section~\ref{sec:rz}, whether $\SoS_\C(Q_M) \le \PT_\C(M)$ holds for \emph{fully symmetric} matrices $M$ remains open.
\item Our citation of \cite{RZ25} refers to v1 (arXiv:2512.01227, December 2025).
\end{enumerate}

\begin{thebibliography}{9}

\bibitem{HWY11}
P.~Hrube\v{s}, A.~Wigderson, and A.~Yehudayoff,
\emph{Non-commutative circuits and the sum-of-squares problem},
Journal of the American Mathematical Society \textbf{24} (2011), no.~3, 871--898.

\bibitem{RZ25}
B.~Rossman and D.~Zhu,
\emph{Multiquadratic sum-of-squares lower bounds imply $\mathrm{VNC}^1 \neq \mathrm{VNP}$},
arXiv:2512.01227v1, December 2025.

\bibitem{Sie26}
J.~P.~Siemaszko,
\emph{On the Hurwitz sum-of-squares problem},
Zenodo, DOI: \texttt{10.5281/zenodo.24024689}, November 2026.

\end{thebibliography}

\end{document}
